\documentclass[11pt, a4paper]{article}

\usepackage[margin=1in]{geometry}
\usepackage{amsmath}
\usepackage{amssymb}
\usepackage{booktabs}
\usepackage{hyperref}
\usepackage{xcolor}
\usepackage{parskip}
\usepackage{microtype}
\usepackage{listings}

\hypersetup{
    colorlinks=true,
    linkcolor=black,
    urlcolor=blue,
    citecolor=black
}

\title{
    \textbf{Phase 2: Numerical Behaviour Audit} \\
    \large Benchmarking Quantum-Inspired Feature Encodings as Explicit
    High-Dimensional Representations in Classical Machine Learning
}
\author{QIE Research Team}
\date{April 2026}

% -------------------------------------------------------------------------
\begin{document}
% -------------------------------------------------------------------------

\maketitle
\tableofcontents
\newpage

% -------------------------------------------------------------------------
\section{Purpose and Scope}
% -------------------------------------------------------------------------

This report documents the Phase~2 numerical behaviour audit for the three
quantum-inspired feature encodings under study: amplitude encoding, angle
encoding, and basis encoding.  The audit is a mandatory pre-training gate.
No performance claims, accuracy results, or downstream model comparisons
appear in this document.

The central questions are:

\begin{enumerate}
    \item Do the encodings satisfy their structural invariants on real data?
    \item What is the geometry of each encoded feature matrix, measured through
          singular value distribution, effective rank, and condition number?
    \item How stable is each encoding under small input perturbations?
\end{enumerate}

All results were produced from the UCI Wine dataset (142 training samples,
13 features, 3 classes) using seed 42.  The audit script is located at
\texttt{src/qie\_research/analysis/numerical\_audit.py} and the raw numbers
are in \texttt{results/metrics/phase\_2\_numerical\_audit.json}.

% -------------------------------------------------------------------------
\section{Encoding Definitions}
% -------------------------------------------------------------------------

Each encoding maps a classical input vector $\mathbf{x} \in \mathbb{R}^d$
to an explicit feature vector $\varphi(\mathbf{x})$.  The three encodings
and their output dimensions are:

\begin{table}[h]
\centering
\begin{tabular}{llll}
\toprule
\textbf{Encoding} & \textbf{Map} & \textbf{Output dim} & \textbf{Differentiable} \\
\midrule
Amplitude & $\varphi(\mathbf{x}) = \hat{\mathbf{x}}$, padded to $2^n$ & $2^{\lceil \log_2 d \rceil}$ & Yes \\
Angle     & $[\cos(\pi x_i/2),\, \sin(\pi x_i/2)]_{i=1}^{d}$, interleaved & $2d$ & Yes \\
Basis     & $b$-bit binary of each quantised $x_i$, concatenated & $d \cdot b$ & No \\
\bottomrule
\end{tabular}
\caption{Summary of the three encodings and their classical feature maps.}
\end{table}

For $d = 13$ input features (UCI Wine) and $b = 8$ bits, the output
dimensions are 16, 26, and 104 respectively.

% -------------------------------------------------------------------------
\section{Formulae}
% -------------------------------------------------------------------------

\subsection{Effective Rank}

Effective rank quantifies how evenly a representation uses its feature
dimensions, via the entropy of the normalised singular value spectrum:

\begin{equation}
    \mathrm{erank}(\mathbf{X}) = \exp\!\left( -\sum_{i} p_i \log p_i \right),
    \qquad p_i = \frac{\sigma_i}{\sum_j \sigma_j}
    \label{eq:erank}
\end{equation}

A rank-1 matrix has $\mathrm{erank} = 1$.  A matrix with a perfectly flat
spectrum has $\mathrm{erank} = d_{\mathrm{out}}$.  Higher effective rank
indicates broader use of representational capacity, though it is not
sufficient on its own to claim practical advantage.

\subsection{Condition Number}

The condition number measures numerical stability as the ratio of the
largest to the smallest non-zero singular value:

\begin{equation}
    \kappa = \frac{\sigma_{\max}}{\sigma_{\min}}
    \label{eq:kappa}
\end{equation}

A high condition number indicates that the feature matrix is nearly
rank-deficient and that small input perturbations can produce large
changes in the encoded space, potentially destabilising optimisation.

\subsection{Noise Stability}

Noise stability measures the mean per-sample Frobenius distance between
the clean and perturbed encodings under additive Gaussian noise
$\varepsilon \sim \mathcal{N}(0,\, \sigma^2)$:

\begin{equation}
    \delta = \frac{\|\varphi(\mathbf{X} + \boldsymbol{\varepsilon}) -
    \varphi(\mathbf{X})\|_F}{n}
    \label{eq:noise}
\end{equation}

All noise stability results in this report use $\sigma = 0.01$.

% -------------------------------------------------------------------------
\section{Invariant Verification}
% -------------------------------------------------------------------------

Each encoding has a structural invariant that must hold for every encoded
sample.  These were verified on all 142 training samples.

\begin{table}[h]
\centering
\begin{tabular}{lll}
\toprule
\textbf{Encoding} & \textbf{Invariant} & \textbf{Result} \\
\midrule
Amplitude & $\sum_i \varphi(\mathbf{x})_i^2 = 1$ & \textcolor{green!60!black}{PASS} \\
Angle     & $\cos^2(\theta_i/2) + \sin^2(\theta_i/2) = 1\ \forall\, i$ & \textcolor{green!60!black}{PASS} \\
Basis     & All values $\in \{0, 1\}$ & \textcolor{green!60!black}{PASS} \\
\bottomrule
\end{tabular}
\caption{Structural invariant verification across all 142 training samples.}
\end{table}

All three encodings satisfy their structural invariants on real data.

% -------------------------------------------------------------------------
\section{Singular Value Distributions}
% -------------------------------------------------------------------------

The top-10 singular values of each encoded feature matrix are reported
below.  These reveal whether the representation distributes energy across
many directions or concentrates it in a few.

\subsection{Amplitude Encoding}

\begin{table}[h]
\centering
\begin{tabular}{cc}
\toprule
\textbf{Rank} & \textbf{Singular value} \\
\midrule
1  & 11.898767 \\
2  & 0.639309  \\
3  & 0.082043  \\
4  & 0.046871  \\
5  & 0.026894  \\
6  & 0.023239  \\
7  & 0.017234  \\
8  & 0.007243  \\
9  & 0.006279  \\
10 & 0.005238  \\
\bottomrule
\end{tabular}
\caption{Top-10 singular values for amplitude encoding ($d_\mathrm{out} = 16$).}
\end{table}

The spectrum is strongly dominated by the first singular value
($\sigma_1 = 11.90$), with a steep drop to $\sigma_2 = 0.64$.  This
indicates that the amplitude-encoded feature matrix is nearly rank-1 ---
almost all variance is concentrated in a single direction.  This is a
direct consequence of L2 normalisation: after normalisation, all samples
lie on the unit hypersphere, and the principal direction captures the
dominant orientation of that cloud.

\subsection{Angle Encoding}

\begin{table}[h]
\centering
\begin{tabular}{cc}
\toprule
\textbf{Rank} & \textbf{Singular value} \\
\midrule
1  & 20.001993 \\
2  & 16.531633 \\
3  & 12.522490 \\
4  & 9.731345  \\
5  & 9.355629  \\
6  & 9.072004  \\
7  & 8.786925  \\
8  & 8.138566  \\
9  & 8.023533  \\
10 & 7.630548  \\
\bottomrule
\end{tabular}
\caption{Top-10 singular values for angle encoding ($d_\mathrm{out} = 26$).}
\end{table}

The spectrum decays gradually with no dominant spike, indicating that the
angle encoding distributes energy broadly across many directions.  The
smallest singular value is 2.65, giving a well-conditioned matrix.

\subsection{Basis Encoding}

\begin{table}[h]
\centering
\begin{tabular}{cc}
\toprule
\textbf{Rank} & \textbf{Singular value} \\
\midrule
1  & 59.763547 \\
2  & 13.209988 \\
3  & 11.430762 \\
4  & 10.954934 \\
5  & 10.793405 \\
6  & 10.290472 \\
7  & 9.907706  \\
8  & 9.418556  \\
9  & 9.333274  \\
10 & 9.260879  \\
\bottomrule
\end{tabular}
\caption{Top-10 singular values for basis encoding ($d_\mathrm{out} = 104$).}
\end{table}

Basis encoding has a large first singular value relative to the rest,
reflecting the large constant offset introduced by the binary encoding of
each feature.  Beyond the first component, the spectrum is relatively
flat, contributing to a high effective rank.

% -------------------------------------------------------------------------
\section{Effective Rank and Condition Number}
% -------------------------------------------------------------------------

\begin{table}[h]
\centering
\begin{tabular}{lcccc}
\toprule
\textbf{Encoding} & \textbf{$d_\mathrm{out}$} & \textbf{erank} &
\textbf{erank / $d_\mathrm{out}$} & \textbf{$\kappa$} \\
\midrule
Amplitude & 16  & 1.3751  & 0.086 & 6650.85 \\
Angle     & 26  & 23.2806 & 0.896 & 7.55    \\
Basis     & 104 & 77.0907 & 0.741 & 108.73  \\
\bottomrule
\end{tabular}
\caption{Effective rank, relative effective rank, and condition number for
each encoding on the UCI Wine training split.}
\end{table}

The results reveal a large spread in representational geometry:

\begin{itemize}
    \item \textbf{Amplitude encoding} has an effective rank of 1.38 out of
    a possible 16, meaning it effectively uses only about 1.4 of its 16
    output dimensions.  The condition number of 6650.85 is very high,
    indicating severe numerical ill-conditioning.  This is consistent with
    the near-rank-1 singular value spectrum observed above and is a
    direct consequence of L2 normalisation collapsing the feature space.

    \item \textbf{Angle encoding} has an effective rank of 23.28 out of 26,
    a relative utilisation of 89.6\%.  The condition number of 7.55 is low,
    indicating a well-conditioned feature matrix.  Angle encoding makes
    near-full use of its output dimensions.

    \item \textbf{Basis encoding} has an effective rank of 77.09 out of 104,
    a relative utilisation of 74.1\%.  The condition number of 108.73 is
    moderate.  The large output dimensionality contributes to a high absolute
    effective rank, though relative utilisation is lower than angle encoding.
\end{itemize}

% -------------------------------------------------------------------------
\section{Noise Stability}
% -------------------------------------------------------------------------

\begin{table}[h]
\centering
\begin{tabular}{lccc}
\toprule
\textbf{Encoding} & \textbf{Noise $\sigma$} &
\textbf{$\delta$ (mean Frobenius / sample)} & \textbf{Relative to $d_\mathrm{out}$} \\
\midrule
Amplitude & 0.01 & $5.0 \times 10^{-6}$ & $3.1 \times 10^{-7}$ \\
Angle     & 0.01 & $1.4 \times 10^{-2}$ & $5.5 \times 10^{-4}$ \\
Basis     & 0.01 & $2.9 \times 10^{-1}$ & $2.8 \times 10^{-3}$ \\
\bottomrule
\end{tabular}
\caption{Noise stability under $\varepsilon \sim \mathcal{N}(0,\, 0.01^2)$.
$\delta$ is the mean per-sample Frobenius distance between clean and noisy
encodings (Equation~\ref{eq:noise}).}
\end{table}

\begin{itemize}
    \item \textbf{Amplitude encoding} is highly stable under noise
    ($\delta = 5.0 \times 10^{-6}$).  This is because L2 normalisation
    absorbs small additive perturbations --- a small $\varepsilon$ changes
    the direction of $\mathbf{x}$ only slightly, so the normalised output
    is nearly unchanged.

    \item \textbf{Angle encoding} shows moderate sensitivity
    ($\delta = 1.4 \times 10^{-2}$).  Small input changes translate to
    small angle changes, which propagate linearly through the cosine and
    sine maps.

    \item \textbf{Basis encoding} is the most sensitive to noise
    ($\delta = 2.9 \times 10^{-1}$).  The quantisation step amplifies
    small perturbations near bin boundaries --- a noise value that pushes
    an input across a quantisation threshold flips a bit, changing the
    encoded output discontinuously.  This is an expected and structural
    property of basis encoding, not an implementation artefact.
\end{itemize}

% -------------------------------------------------------------------------
\section{Summary and Readiness Assessment}
% -------------------------------------------------------------------------

\begin{table}[h]
\centering
\begin{tabular}{lcccc}
\toprule
\textbf{Encoding} & \textbf{Invariant} & \textbf{erank} &
\textbf{$\kappa$} & \textbf{Noise $\delta$} \\
\midrule
Amplitude & \textcolor{green!60!black}{PASS} & 1.38  & 6650.85 & $5.0 \times 10^{-6}$ \\
Angle     & \textcolor{green!60!black}{PASS} & 23.28 & 7.55    & $1.4 \times 10^{-2}$ \\
Basis     & \textcolor{green!60!black}{PASS} & 77.09 & 108.73  & $2.9 \times 10^{-1}$ \\
\bottomrule
\end{tabular}
\caption{Summary of Phase 2 numerical audit results.}
\end{table}

All three encodings pass their structural invariant checks and are
numerically well-defined on the UCI Wine dataset.  The audit surfaces
three properties that are directly relevant to the Phase~4 and Phase~5
analyses:

\begin{enumerate}
    \item Amplitude encoding is nearly rank-1 and severely ill-conditioned.
    Any downstream model operating on amplitude-encoded features is working
    in a near-one-dimensional space.  This will need to be reconciled with
    its Phase~4 performance results and explained through the conditioning
    and NTK analyses in Phase~5.

    \item Angle encoding has near-full effective rank utilisation and a low
    condition number.  It is the most geometrically well-behaved encoding
    of the three under these conditions.

    \item Basis encoding is sensitive to noise due to the discontinuous
    quantisation step.  This is a structural limitation that must be
    reported alongside any Phase~4 performance results and must be
    highlighted in the practical overhead discussion.
\end{enumerate}

\textbf{Phase 3 (Dataset Freeze) may begin.}  No performance claims have
been made in this report.

% -------------------------------------------------------------------------
\section{Artifact References}
% -------------------------------------------------------------------------

\begin{itemize}
    \item Audit script:
    \texttt{src/qie\_research/analysis/numerical\_audit.py}
    \item Raw audit output:
    \texttt{results/metrics/phase\_2\_numerical\_audit.json}
    \item Phase record:
    \texttt{records/phases/phase-2-numerical-behavior.md}
    \item Encoding implementations:
    \texttt{src/qie\_research/encodings/}
\end{itemize}

\end{document}
